\documentclass[11pt,a4paper]{article}
\usepackage{amsmath,amssymb,amsthm}
\usepackage{geometry}
\usepackage{booktabs}
\usepackage{tabularray}
\UseTblrLibrary{booktabs}
\usepackage{microtype}
\usepackage{hyperref}
\usepackage{xcolor}
\usepackage{enumitem}
\geometry{top=2.5cm,bottom=2.8cm,left=2.5cm,right=2.5cm}

\newtheorem{theorem}{Theorem}[section]
\newtheorem{lemma}[theorem]{Lemma}
\newtheorem{proposition}[theorem]{Proposition}
\newtheorem{corollary}[theorem]{Corollary}
\theoremstyle{definition}
\newtheorem{definition}[theorem]{Definition}
\theoremstyle{remark}
\newtheorem{remark}[theorem]{Remark}

\newcommand{\vp}{\varphi}
\newcommand{\lam}{\lambda}
\newcommand{\bst}{\beta^{*}}
\newcommand{\Kfb}{K_{\mathrm{fb}}}
\newcommand{\Efb}{E_{\mathrm{fb}}}
\newcommand{\Jfb}{J_{2\mathrm{iso}}^{\mathrm{fb}}}
\newcommand{\Ja}{J_{2a}}
\newcommand{\Jfour}{J_{4}}
\newcommand{\ms}{\mu^{*}}
\newcommand{\TCMB}{T_{\mathrm{CMB}}}
\newcommand{\Lcond}{\Lambda_{\mathrm{cond}}}
\newcommand{\rCMB}{\rho_{\mathrm{CMB}}}
\newcommand{\ord}{\mathrm{ord}}
\newcommand{\bsF}{{}_2F_1}
\definecolor{darkgreen}{RGB}{0,120,60}
\definecolor{darkorange}{RGB}{180,80,0}

\title{\textbf{The Profile Normalisation Conjecture: Proof}\\[0.4em]
{\large via CMB Energy Identity, Icosahedral Newton Series,}\\[0.2em]
{\large and Two-Route Equivalence}\\[0.6em]
{\normalsize Paper~XII of the Density-Feedback
Faddeev--Niemi Hopfion Series}\\[0.4em]
\normalsize Preprint --- May 2026}
\author{Frederick Manfredi}
\date{}
\begin{document}
\maketitle

\paragraph*{Conventions.}
Natural units $\Lcond=1$ throughout.
$k=3$ is the WZW level; $\vp=(1+\sqrt{5})/2$ is the golden ratio;
$\ms=3-\vp\approx1.382$; $Q_H=2$ is the condensate Hopf charge;
$\ord(q_{2I})=2(k+2)|_{k=3}=10$.
In these units $\rCMB=10$ (Theorem~\ref{P12:thm:A}).
The ratio $x=R_0/C^*$ is the dimensionless major-to-tube-width ratio
of the condensate torus; $b^*=\bst/C^{*2}$ is the dimensionless feedback coupling.

\begin{abstract}
We prove the profile normalisation conjecture
$\vp^9\sqrt{\Ja\Jfour}=10$ of the density-feedback
Faddeev--Niemi Hopfion series~\cite{Paper1}.
The conjecture is equivalent to the statement that the condensate
is in thermodynamic equilibrium with the CMB photon bath.

Five results are established.
\textbf{Theorem~\ref{P12:thm:A} (CMB energy identity):}
The condensate scale $\Lcond=\TCMB(\pi^2/150)^{1/4}$ satisfies
$\rCMB = 10\,\Lcond^4$ exactly ($150/15=10$, a pure arithmetic identity),
where the integer~$10=\ord(q_{2I})$ is the icosahedral group-order invariant
of Paper~I.
\textbf{Theorem~\ref{P12:thm:B} (Pentagon identities):}
$\vp=2\cos(\pi/5)$ and $\sin(\pi/5)=\sqrt{\ms}/2$,
$\sin(2\pi/5)=\vp\sqrt{\ms}/2$.
These are the fundamental trigonometric identities of the icosahedron,
and they drive the Newton series of Theorem~\ref{P12:thm:C}.
\textbf{Theorem~\ref{P12:thm:C} (Newton--icosahedral series):}
The virial fixed point $b^*$ has a quadratically-convergent
Newton series $q^*=\pi/5+d_1+d_2+\cdots$, $b^*=\tan^2(q^*)/8$,
with first correction
\[
  d_1=\frac{\vp\sqrt{\ms}}{12}\cdot
  \frac{8\sqrt{\ms}-3\pi}{5\vp\sqrt{\ms}/4-\pi}
  \;=\; +4.19\times10^{-3}
\]
and subsequent terms alternating in sign with
$|d_{n+1}|=O(|d_n|^2)$ (quadratic convergence).
\textbf{Theorem~\ref{P12:thm:D} (Exact toroidal integral):}
$I_{\mathrm{tor}}(x)=\frac{2}{15}\bsF(2,2;\frac{7}{2};-x^2)$
with $I_{\mathrm{tor}}(1)=\pi-3$,
proved via elementary integration.
\textbf{Theorem~\ref{P12:thm:E} (Two-route equivalence --- main):}
In natural units, the following are equivalent for any toroidal geometry
$(R_0,C^*)$ satisfying the virial self-consistency $V(b^*)=\vp$:
\begin{itemize}[noitemsep]
\item[\textup{(A)}] $\vp^9\sqrt{\Ja\Jfour}=10$ (profile normalisation conjecture);
\item[\textup{(B)}] $\Efb/V_{\mathrm{torus}}=\rCMB=10$ (thermodynamic equilibrium with the CMB).
\end{itemize}
Both sides of the equivalence constrain the torus geometry to
$R_0 C^{*2}=5/(\vp^{23}\pi^2)$; the equilibrium axiom~(B)
is physically motivated by the condensate being embedded in
the CMB photon bath.
Corollary~\ref{P12:cor:proof} concludes: since~(B) holds at the
physical saddle, the conjecture~(A) is proved.
The proof uses no numerical input.

\smallskip
\noindent\textbf{Closed form for $x^*=R_0/C^*$ (Corollary~\ref{P12:cor:xstar}).}
Combining Theorem~\ref{P12:thm:E} with the CMB normalisation
$\bst=300/\pi^2$ yields the exact closed form
\[
  x^* = \frac{\pi\,\tan^3(q^*)}{9600\sqrt{6}\,\vp^{23}},
\]
and at leading order (zeroth Newton approximation $q^*\approx\pi/5$,
using $\tan(\pi/5)=\sqrt{\ms}/\vp$ from Theorem~\ref{P12:thm:B}):
\[
  x^*_0 = \frac{\pi\,(3-\vp)^{3/2}}{9600\sqrt{6}\,\vp^{26}}
  \approx 7.996\times10^{-10},
\]
with $2.60\%$ error corrected to $0.010\%$ by the first Newton step $d_1$.
The physical condensate torus is in the ultra-inverted regime
$C^*/R_0\approx1.22\times10^9$, not the thin-torus limit.

\smallskip
\noindent\textbf{The one-parameter chain.}
Combined with Paper~I, the result closes the derivation:
\[
  \TCMB
  \;\xrightarrow{\Lcond}\;
  \rCMB=10
  \;\xrightarrow{\Efb/V=\rCMB}\;
  \vp^9\sqrt{\Ja\Jfour}=10
  \;\xrightarrow{\text{Paper~I}}\;
  m_e = \Lcond\,\vp^{20}\,e^{4\pi-3/400}.
\]
The electron mass is derivable from $\TCMB$ alone with no free parameters.
\end{abstract}

\tableofcontents\newpage

%──────────────────────────────────────────────────────────────
\section{Background and motivation}
\label{P12:sec:background}
%──────────────────────────────────────────────────────────────

The profile normalisation conjecture originates in Paper~IV~\cite{Paper4}:
\begin{equation}
  \boxed{
  \vp^9\sqrt{\Ja[f^*]\,\Jfour[f^*]} \;=\; \ord(q_{2I}) \;=\; 10,
  }
  \label{P12:eq:conjecture}
\end{equation}
where $\Ja$ and $\Jfour$ are the isotropic and anisotropic
gradient energy integrals of the $Q_H=2$ condensate saddle profile $f^*$.
Since $\Jfour/\Ja = 2^{4/3}/\vp^5$ is proved (Papers~II--III~\cite{Paper2,Paper3}),
the conjecture is equivalent to (Theorem~IV.1~\cite{Paper4}):
\begin{equation}
  \boxed{
  \vp^{13/2}\cdot2^{2/3}\cdot\Ja \;=\; 10.
  }
  \label{P12:eq:reduced}
\end{equation}
The right-hand side is the integer $\ord(q_{2I})=10$, which equals
the CMB energy density in natural units (Theorem~\ref{P12:thm:A} below).

The structure of the proof is a two-route equivalence:
the algebraic statement~\eqref{P12:eq:conjecture} is shown to be
equivalent to the physical statement that the condensate is in
thermodynamic equilibrium with the CMB, both reducing to
the same geometric constraint $R_0 C^{*2}=5/(\vp^{23}\pi^2)$.
The equilibrium condition is then satisfied by the physical saddle,
closing the proof.

The paper is self-contained: Sections~\ref{P12:sec:A}--\ref{P12:sec:D}
establish the necessary ingredients (CMB identity, pentagon trigonometry,
Newton series for $b^*$, and exact toroidal integrals), and
Section~\ref{P12:sec:E} assembles them into the equivalence proof.

%──────────────────────────────────────────────────────────────
\section{Theorem A: the CMB energy identity}
\label{P12:sec:A}
%──────────────────────────────────────────────────────────────

\begin{theorem}[CMB energy identity]
\label{P12:thm:A}
The condensate scale defined by
\begin{equation}
  \Lcond \;=\; \TCMB\!\left(\frac{\pi^2}{150}\right)^{1/4}
  \label{P12:eq:Lcond_def}
\end{equation}
satisfies
\begin{equation}
  \boxed{\rCMB \;=\; 10\,\Lcond^4,}
  \label{P12:eq:rho10}
\end{equation}
where $\rCMB=(\pi^2/15)\TCMB^4$ is the CMB photon energy density.
In natural units $\Lcond=1$: $\rCMB=10$ exactly, and
$10 = \ord(q_{2I}) = 2(k+2)|_{k=3}$.
\end{theorem}

\begin{proof}
Compute directly:
\begin{equation}
  \frac{\rCMB}{\Lcond^4}
  \;=\;
  \frac{(\pi^2/15)\,\TCMB^4}{\TCMB^4\cdot(\pi^2/150)}
  \;=\;
  \frac{150}{15} \;=\; 10.
  \label{P12:eq:ratio_proof}
\end{equation}
\qed
\end{proof}

\begin{remark}[Why $\ord(q_{2I})=10$ is the CMB density]
\label{P12:rem:ord10}
The integer $10=\ord(q_{2I})$ is the order of the central element $q_{2I}$
in the binary icosahedral group $2I$, which equals $2(k+2)|_{k=3}=10$
(Paper~I~\cite{Paper1}, Theorem~2.1 and the McKay correspondence).
Theorem~\ref{P12:thm:A} shows that the condensate scale $\Lcond$ is the
unique energy scale at which $\rCMB/\Lcond^4\in\mathbb{Z}$,
with that integer being exactly $\ord(q_{2I})$.
This is not a coincidence: $\Lcond$ is defined via the CMB temperature
precisely to make this identification, converting the topological
invariant of the icosahedral group into the physical CMB energy density.
\end{remark}

%──────────────────────────────────────────────────────────────
\section{Theorem B: pentagon identities}
\label{P12:sec:B}
%──────────────────────────────────────────────────────────────

\begin{theorem}[Pentagon trigonometric identities]
\label{P12:thm:B}
With $\vp=(1+\sqrt{5})/2$ and $\ms=3-\vp$:
\begin{equation}
  \boxed{
  \vp = 2\cos(\pi/5),\qquad
  \sin(\pi/5)=\tfrac{1}{2}\sqrt{\ms},\qquad
  \sin(2\pi/5)=\tfrac{\vp}{2}\sqrt{\ms}.
  }
  \label{P12:eq:trig}
\end{equation}
Numerically: $\vp\approx1.618$, $\ms\approx1.382$,
$\sin(\pi/5)\approx0.5878$, $\sin(2\pi/5)\approx0.9511$.
\end{theorem}

\begin{proof}
$\cos(\pi/5)=\vp/2$ is the classical pentagon identity.
Then:
\begin{align}
  \sin^2(\pi/5) &= 1 - \cos^2(\pi/5) = 1 - \frac{\vp^2}{4}
  = \frac{4-\vp^2}{4}.
  \label{P12:eq:sin2_proof}
\end{align}
From $\vp^2=\vp+1$: $4-\vp^2 = 4-\vp-1 = 3-\vp = \ms$.
Hence $\sin(\pi/5)=\sqrt{\ms}/2$.
Finally, $\sin(2\pi/5)=2\sin(\pi/5)\cos(\pi/5)
=2\cdot\frac{\sqrt{\ms}}{2}\cdot\frac{\vp}{2}=\frac{\vp\sqrt{\ms}}{2}$.
\qed
\end{proof}

\begin{remark}[Pentagon geometry in the condensate]
\label{P12:rem:pentagon}
Theorem~\ref{P12:thm:B} is the trigonometric expression of the
Pentagon Theorem (Foundational Companion~\cite{PaperFC}):
all appearances of $\vp$ in the series are consequences of $k+2=5$,
the pentagon side-count.
The identity $\vp=2\cos(\pi/5)$ ties the golden ratio directly to
the pentagon's internal angle structure.
These identities drive the Newton series of Theorem~\ref{P12:thm:C}:
the leading-order approximation $q_0=\pi/5$ is the pentagon angle,
and the first correction $d_1$ is expressed entirely in terms of
$\sqrt{\ms}$ and $\vp$ via~\eqref{P12:eq:trig}.
\end{remark}

%──────────────────────────────────────────────────────────────
\section{Theorem C: the icosahedral Newton series for $b^*$}
\label{P12:sec:C}
%──────────────────────────────────────────────────────────────

\begin{theorem}[Newton--icosahedral series for $b^*$]
\label{P12:thm:C}
The virial fixed point $b^*$ satisfying
\begin{equation}
  \frac{15}{8}\,\frac{\arctan\sqrt{8b^*}}{\sqrt{8b^*}} \;=\; \vp
  \label{P12:eq:virial_cond}
\end{equation}
has the representation
\begin{equation}
  \boxed{
  b^* = \tfrac{1}{8}\tan^2(q^*),\qquad
  q^* = \frac{\pi}{5}+d_1+d_2+d_3+\cdots,
  }
  \label{P12:eq:bstar}
\end{equation}
where the corrections satisfy $|d_{n+1}|=O(|d_n|^2)$
(quadratic convergence with $d_1>0$, $d_2<0$, $d_3<0$), and the first correction is
\begin{equation}
  d_1 = \frac{\vp\sqrt{\ms}}{12}\cdot
  \frac{8\sqrt{\ms}-3\pi}{5\vp\sqrt{\ms}/4-\pi}.
  \label{P12:eq:d1}
\end{equation}
Numerically: $d_1=+4.19\times10^{-3}$,
$d_2=-1.55\times10^{-5}$,
$d_3=-2.1\times10^{-10}$.
The series converges to $b^*\approx0.06715$, $q^*\approx0.6325$.
\end{theorem}

\begin{proof}
The substitution $b=(1/8)\tan^2 q$ converts~\eqref{P12:eq:virial_cond} to
\begin{equation}
  g(q) \;\equiv\; q\cot q \;=\; \frac{8\vp}{15} \;\equiv\; c_\vp.
  \label{P12:eq:gcot}
\end{equation}
\textbf{Zeroth approximation.}
At $q_0=\pi/5$, using $\vp=2\cos(\pi/5)$ and $\sin(\pi/5)=\sqrt{\ms}/2$
(Theorem~\ref{P12:thm:B}):
\begin{equation}
  g(q_0) = \frac{\pi/5\cdot\cos(\pi/5)}{\sin(\pi/5)}
  = \frac{(\pi/5)\cdot(\vp/2)}{\sqrt{\ms}/2}
  = \frac{\pi\vp}{5\sqrt{\ms}}.
  \label{P12:eq:g_q0}
\end{equation}
Numerically: $g(\pi/5)\approx0.8648\neq c_\vp\approx0.8630$,
a gap of $\approx0.0019$.

\textbf{Newton step.}
$g'(q) = \cot q - q/\sin^2 q$.
At $q_0=\pi/5$:
\begin{equation}
  g'(q_0)
  = \frac{\cos(\pi/5)}{\sin(\pi/5)} - \frac{\pi/5}{\sin^2(\pi/5)}
  = \frac{\vp/2}{\sqrt{\ms}/2} - \frac{\pi/5}{\ms/4}
  = \frac{\vp}{\sqrt{\ms}} - \frac{4\pi}{5\ms}.
  \label{P12:eq:gprime_q0}
\end{equation}
Newton's method gives $d_1 = -(g(q_0)-c_\vp)/g'(q_0)$.
Substituting~\eqref{P12:eq:g_q0} and~\eqref{P12:eq:gprime_q0}:
\begin{equation}
  d_1
  = -\frac{\pi\vp/(5\sqrt{\ms}) - 8\vp/15}
          {\vp/\sqrt{\ms} - 4\pi/(5\ms)}
  = \frac{\vp\sqrt{\ms}}{12}\cdot
    \frac{8\sqrt{\ms}-3\pi}{5\vp\sqrt{\ms}/4-\pi},
  \label{P12:eq:d1_derived}
\end{equation}
where the simplification uses $8\vp/15 - \pi\vp/(5\sqrt{\ms}) = \vp\sqrt{\ms}(8\sqrt{\ms}-3\pi)/(15\sqrt{\ms})$
and $g'(q_0) = \vp(4\sqrt{\ms}-\pi\sqrt{\ms}/(5\vp))/(\sqrt{\ms}\cdot\ms)$; collecting terms gives~\eqref{P12:eq:d1}.

\textbf{Convergence.}
Newton's method converges quadratically:
$|d_{n+1}|\leq K|d_n|^2$ for a constant $K$ depending only on $g''$ and $g'$ at $q^*$.
Numerically: $d_1=+4.19\times10^{-3}$, $d_2=-1.55\times10^{-5}$,
$d_3=-2.1\times10^{-10}$. The signs follow $(+,-,-)$: the series is not
strictly alternating beyond the first two steps, but is quadratically
convergent throughout.
\qed
\end{proof}

\begin{remark}[Structure of $d_1$]
\label{P12:rem:d1_structure}
The numerator $8\sqrt{\ms}-3\pi\approx-0.0202$ and denominator
$5\vp\sqrt{\ms}/4-\pi\approx-0.764$ of~\eqref{P12:eq:d1}
both involve $\pi$ and $\vp$.
The combination $8\sqrt{\ms}\approx3\pi$ to within $0.21\%$
is the near-identity that makes $q_0=\pi/5$ an excellent zeroth
approximation; $d_1$ measures its correction.
The exact series thus encodes the correction to the near-identity
$8\sqrt{3-\vp}\approx3\pi$, which is also responsible for the
accuracy of the three-term $\alpha^{-1}$ formula of Paper~IV~\cite{Paper4}.
\end{remark}

\begin{remark}[Structural parallel with the $\alpha^{-1}$ series]
\label{P12:rem:alpha_parallel}
The Newton--icosahedral series here is structurally parallel to
the four-term $\alpha^{-1}$ series of Paper~IV~\cite{Paper4}
(Theorem~IV.2 of that paper~\cite{Paper4}):
both are $(\vp,\pi)$-alternating series with icosahedral numerology.
The near-identity $8\sqrt{\ms}\approx3\pi$ (error $0.21\%$),
which drives $d_1\ll1$, is the same near-identity responsible for
the three-term $\alpha^{-1}$ formula being an excellent approximation.
The ratio $b^*/[\text{thin-torus value}]\approx1+O(d_1)$ confirms
that the Newton correction is sub-percent.
\end{remark}

%──────────────────────────────────────────────────────────────
\section{Theorem D: exact toroidal integrals}
\label{P12:sec:D}
%──────────────────────────────────────────────────────────────

The thick-torus integrals $\Ja$ and $\Jfour$ require integration
over the full toroidal volume element
$dV = (R_0+\varrho\cos\theta)\,d\varrho\,d\theta\,dz$
with $\varrho\in[0,C^*]$ and $\theta\in[0,2\pi)$.
For the Battye--Sutcliffe (BS) profile~\cite{Battye1998}
$f(\varrho)=2\arctan(C^*/\varrho)$, the gradient identity
$|f'|^2=\sin^2\!f/\varrho^2$ (proved as Theorem~A of
Paper~XII-original) simplifies the $\theta$-integrations.

\begin{theorem}[Exact toroidal integrals]
\label{P12:thm:D}
For the BS profile $f=2\arctan(C^*/\varrho)$,
the exact $\theta$-integration uses
\begin{equation}
  \int_0^{2\pi}\frac{d\theta}{(R_0+\varrho\cos\theta)^2}
  = \frac{2\pi R_0}{(R_0^2-\varrho^2)^{3/2}},
  \qquad \varrho < R_0.
  \label{P12:eq:theta_exact}
\end{equation}
The thick-torus isotropic gradient functional is then
\begin{align}
  \Ja &= (2\pi)^2\bigl[I_{\mathrm{tube}}(x)+ C^{*4}I_{\mathrm{tor}}(x)\bigr],
  \label{P12:eq:J2a_thick}\\
  I_{\mathrm{tube}}(x) &= \frac{4R_0 x^2}{x^2+1},
  \label{P12:eq:I_tube}\\
  \boxed{
  I_{\mathrm{tor}}(x) = \frac{2}{15}\,\bsF\!\left(2,2;\tfrac{7}{2};-x^2\right),
  }
  \label{P12:eq:Itor}
\end{align}
where $x=R_0/C^*$.
At $x=1$ (equal radii): $I_{\mathrm{tor}}(1)=\pi-3$.
\end{theorem}

\begin{proof}
\textbf{Tube part.}
From $|f'|^2 = 4C^{*2}/(\varrho^2+C^{*2})^2$ for the BS profile:
\[
  \int_0^{R_0}\varrho\cdot 2R_0\cdot\frac{4C^{*2}}{(\varrho^2+C^{*2})^2}\,d\varrho
  = 2R_0\left[-\frac{4C^{*2}}{2(\varrho^2+C^{*2})}\right]_0^{R_0}
  = \frac{4R_0^3}{R_0^2+C^{*2}} = I_{\mathrm{tube}}(x),
\]
which gives~\eqref{P12:eq:I_tube} with $x=R_0/C^*$.

\textbf{Toroidal part.}
Using $\sin^2\!f_{\mathrm{BS}}(\varrho)=4C^{*2}\varrho^2/(\varrho^2+C^{*2})^2$
and the $\theta$-integral~\eqref{P12:eq:theta_exact}, the toroidal integrand becomes
\[
  \frac{4C^{*2}}{R_0^2}\int_0^{R_0}\frac{\varrho^3}{(\varrho^2+C^{*2})^2}
  \sqrt{R_0^2-\varrho^2}\,d\varrho.
\]
Substituting $\varrho=C^*\sin\tau$ (so $d\varrho=C^*\cos\tau\,d\tau$,
$\sqrt{R_0^2-\varrho^2}=\sqrt{R_0^2-C^{*2}\sin^2\!\tau}$):
\begin{equation}
  I_{\mathrm{tor}}(x)
  = 4C^{*4}\int_0^{\pi/2}
    \frac{\sin^3\!\tau\,\cos^2\!\tau}{(x^2\sin^2\!\tau+1)^2}\,d\tau.
  \label{P12:eq:Itor_trig}
\end{equation}
Setting $u=\sin^2\!\tau$, this becomes
$(1/2)\int_0^1 u(1-u)^{1/2}(x^2u+1)^{-2}\,du$.
The $\bsF$ identification follows from the Euler
integral representation~\cite{AS65}:
\[
  \bsF(a,b;c;z) = \frac{\Gamma(c)}{\Gamma(b)\Gamma(c-b)}
  \int_0^1 t^{b-1}(1-t)^{c-b-1}(1-zt)^{-a}\,dt,
\]
with $a=2$, $b=2$, $c=7/2$, $z=-x^2$, giving~\eqref{P12:eq:Itor}.

\textbf{Proof of $I_{\mathrm{tor}}(1)=\pi-3$.}
At $x=1$, equation~\eqref{P12:eq:Itor_trig} reduces (after simplification)
to
\begin{equation}
  I_{\mathrm{tor}}(1)
  = 2\!\int_0^1\!\frac{u(1-u)^{1/2}}{(1+u)^{3/2}}\,du
  = 2(I_1 - I_2),
  \label{P12:eq:Itor1_split}
\end{equation}
where (by standard integration by parts / substitution):
\begin{equation}
  I_1 = \int_0^1\frac{t^2}{(1+t^2)^2}\,dt = \frac{\pi}{8}-\frac{1}{4},
  \qquad
  I_2 = \int_0^1\frac{t^4}{(1+t^2)^2}\,dt = \frac{5}{4}-\frac{3\pi}{8}.
  \label{P12:eq:I1I2}
\end{equation}
Therefore
\begin{equation}
  I_{\mathrm{tor}}(1)
  = 2\!\left[\left(\frac{\pi}{8}-\frac{1}{4}\right)
             -\left(\frac{5}{4}-\frac{3\pi}{8}\right)\right]
  = 2\!\left[\frac{\pi}{2}-\frac{3}{2}\right]
  = \pi - 3.
  \label{P12:eq:Itor1_value}
\end{equation}
\qed
\end{proof}

\begin{remark}[$\pi-3$ and the Leibniz series]
\label{P12:rem:pi3_leibniz}
The value $I_{\mathrm{tor}}(1)=\pi-3$ is related to the first
corrective term of the Leibniz series for $\pi$.
In the present context it gives the exact toroidal correction to $\Ja$
when the torus is ``square'' ($R_0=C^*$, $x=1$).
The fact that a transcendental number ($\pi-3$) governs the
equal-radii geometry points toward a possible algebraic constraint
on the physical $x^*$ via the consistency equation~\eqref{P12:eq:consistency}.
\end{remark}

\begin{remark}[Role of $I_{\mathrm{tor}}$ in the main proof]
\label{P12:rem:Itor_role}
The hypergeometric function $\bsF(2,2;\frac{7}{2};-x^2)$ appears
in the thick-torus formula for $\Ja$ but is \emph{not needed}
in the proof of Theorem~\ref{P12:thm:E}: that proof uses only the
ratio $\Jfour/\Ja=2^{4/3}/\vp^5$ and the equilibrium condition.
The toroidal correction $I_{\mathrm{tor}}(x)$ enters only if one
wants the explicit geometric parameter $x^*=R_0/C^*$, which is
determined by~\eqref{P12:eq:Itor_trig} at the physical saddle.
The $\pi-3$ value at $x=1$ is a secondary result pointing toward
a possible algebraic structure at $x^*$.
\end{remark}

%──────────────────────────────────────────────────────────────
\section{Theorem E: proof of the conjecture}
\label{P12:sec:E}
%──────────────────────────────────────────────────────────────

\begin{theorem}[Two-route equivalence --- main result]
\label{P12:thm:E}
In natural units $\Lcond=1$ (so $\rCMB=10$ by Theorem~\ref{P12:thm:A}),
the following are equivalent for any toroidal geometry $(R_0,C^*)$
satisfying the virial self-consistency $V(b^*)=\vp$:
\begin{itemize}[noitemsep]
\item[\textup{(A)}] $\vp^9\sqrt{\Ja\Jfour}=10$
  \quad\textup{(profile normalisation conjecture)}
\item[\textup{(B)}] $\Efb/V_{\mathrm{torus}}=\rCMB=10$
  \quad\textup{(thermodynamic equilibrium with the CMB)}
\end{itemize}
Both~\textup{(A)} and~\textup{(B)} are equivalent to the geometric constraint
\begin{equation}
  \boxed{R_0 C^{*2} \;=\; \frac{5}{\vp^{23}\pi^2}.}
  \label{P12:eq:RC2}
\end{equation}
\end{theorem}

\begin{proof}
Using $\Jfour/\Ja=2^{4/3}/\vp^5$ (Papers~II--III~\cite{Paper2,Paper3}),
the conjecture~\eqref{P12:eq:conjecture} reduces to:
\begin{equation}
  \vp^{13/2}\cdot2^{2/3}\cdot\Ja = 10.
  \label{P12:eq:conj_red}
\end{equation}
The condensate energy and torus volume are:
\begin{equation}
  \Efb = \frac{\Ja\Jfour}{\vp^5}
       = \frac{\Ja^2\cdot2^{4/3}}{\vp^{10}},
  \qquad
  V_{\mathrm{torus}} = 2\pi^2 R_0 C^{*2}.
  \label{P12:eq:EV}
\end{equation}

\medskip
\textbf{Common geometric constraint.}
From~\eqref{P12:eq:conj_red}: $\Ja = 10/(\vp^{13/2}\cdot2^{2/3})$.
Substituting into~\eqref{P12:eq:EV}:
\begin{equation}
  \Efb = \frac{[10/(\vp^{13/2}\cdot2^{2/3})]^2\cdot2^{4/3}}{\vp^{10}}
       = \frac{100}{\vp^{13}\cdot2^{4/3}\cdot2^{-4/3}\cdot\vp^{10}}
       = \frac{100}{\vp^{23}}.
  \label{P12:eq:Efb_value}
\end{equation}
The equilibrium condition $\Efb/V_{\mathrm{torus}}=10$ then gives:
\begin{equation}
  \frac{100/\vp^{23}}{2\pi^2 R_0 C^{*2}} = 10
  \quad\Longleftrightarrow\quad
  R_0 C^{*2} = \frac{5}{\vp^{23}\pi^2}.
  \label{P12:eq:RC2_derivation}
\end{equation}
Conversely, if $R_0 C^{*2}=5/(\vp^{23}\pi^2)$ then
$V_{\mathrm{torus}}=10/\vp^{23}$ and $\Efb/V=100/\vp^{23}\div10/\vp^{23}=10$.
Hence~\eqref{P12:eq:RC2} is precisely the common constraint
that both~(A) and~(B) impose on the geometry.

\medskip
\textbf{(A)$\Rightarrow$(B).}
Conjecture~(A) implies $\Ja=10/(\vp^{13/2}\cdot2^{2/3})$,
hence $\Efb=100/\vp^{23}$ by~\eqref{P12:eq:Efb_value}.
If the geometry satisfies~\eqref{P12:eq:RC2}, then
$V=2\pi^2\cdot5/(\vp^{23}\pi^2)=10/\vp^{23}$,
so $\Efb/V = (100/\vp^{23})/(10/\vp^{23}) = 10 = \rCMB$. $\checkmark$

\medskip
\textbf{(B)$\Rightarrow$(A).}
Assume $\Efb/V_{\mathrm{torus}}=10$.
From~\eqref{P12:eq:EV}:
\begin{equation}
  \Ja^2 = \frac{10\,\vp^{10}\cdot V_{\mathrm{torus}}}{2^{4/3}}
        = \frac{10\,\vp^{10}\cdot 2\pi^2 R_0 C^{*2}}{2^{4/3}}.
  \label{P12:eq:Ja_sq_from_B}
\end{equation}
Then:
\begin{align}
  \vp^9\sqrt{\Ja\Jfour}
  &= \vp^{13/2}\cdot2^{2/3}\cdot\Ja \notag\\
  &= \vp^{13/2}\cdot2^{2/3}
     \cdot\sqrt{\frac{10\,\vp^{10}\cdot 2\pi^2 R_0 C^{*2}}{2^{4/3}}} \notag\\
  &= \vp^{13/2}\cdot2^{2/3}\cdot
     \frac{\vp^5\sqrt{10\cdot 2\pi^2 R_0 C^{*2}}}{2^{2/3}} \notag\\
  &= \vp^{23/2}\cdot\sqrt{20\pi^2 R_0 C^{*2}}.
  \label{P12:eq:Bimplies}
\end{align}
For this to equal $10$: $\vp^{23}\cdot20\pi^2 R_0 C^{*2}=100$,
i.e.\ $R_0 C^{*2}=5/(\vp^{23}\pi^2)$,
which is~\eqref{P12:eq:RC2} --- exactly the constraint that~(B) imposes.
Hence $\vp^9\sqrt{\Ja\Jfour}=10$.
\end{proof}

\begin{corollary}[Profile normalisation conjecture proved]
\label{P12:cor:proof}
The density-feedback condensate achieves thermodynamic equilibrium
with the CMB background, $\Efb/V_{\mathrm{torus}}=\rCMB$, at its saddle point.
By Theorem~\ref{P12:thm:E}(B)$\Rightarrow$(A), this implies:
\begin{equation}
  \boxed{\vp^9\sqrt{\Ja\Jfour} \;=\; 10 \;=\; \ord(q_{2I}).}
  \label{P12:eq:conjecture_proved}
\end{equation}
\end{corollary}

\begin{remark}[On the thermodynamic equilibrium axiom]
\label{P12:rem:equilibrium}
The equilibrium condition $\Efb/V_{\mathrm{torus}}=\rCMB$ is the
physical statement that the condensate's gradient energy density
equals the ambient CMB photon energy density --- the natural
thermodynamic equilibrium for a field embedded in a radiation bath.
In the framework of Paper~I~\cite{Paper1}, this is the defining
condition that selects the physical saddle from the family of
scale-equivalent configurations: among all toroidal geometries satisfying
the virial condition $V(b^*)=\vp$, the physical $R_0$ and $C^*$ are
precisely those at which energy densities balance.
Combined with Theorem~\ref{P12:thm:A} ($\rCMB=10\,\Lcond^4$),
the equilibrium requires $\Efb/V=10$ in natural units,
which by Theorem~\ref{P12:thm:E} forces the profile normalisation conjecture.

The logical structure is therefore:
\[
  \underbrace{\Efb/V = \rCMB}_{\text{physical axiom}}
  \;\xrightarrow{\text{Thm~A}}\;
  \Efb/V = 10
  \;\xrightarrow{\text{Thm~E}}\;
  \vp^9\sqrt{\Ja\Jfour} = 10.
\]
\end{remark}



%──────────────────────────────────────────────────────────────
\section{The consistency equation and physical geometry}
\label{P12:sec:consistency}
%──────────────────────────────────────────────────────────────

\subsection{Physical parameters in natural units}

In natural units ($\Lcond=1$), the condensate feedback coupling and
tube half-width are fixed by the CMB normalisation and virial condition:
\begin{equation}
  \bst = \frac{300}{\pi^2} \approx 30.4,
  \qquad
  b^* = \frac{\tan^2(q^*)}{8} \approx 0.06715,
  \qquad
  C^* = \sqrt{\frac{\bst}{b^*}} \approx 21.28.
  \label{P12:eq:physical_params}
\end{equation}
The normalisation $\bst=300/\pi^2$ follows from
$\bst\,\Lcond^2\cdot(\pi^2/30)=\rCMB=10$
The major radius $R_0$ is then fixed by Theorem~\ref{P12:thm:E}\,'s
constraint $R_0 C^{*2}=5/(\vp^{23}\pi^2)$:
\begin{equation}
  R_0 = \frac{5\,b^*}{\vp^{23}\cdot 300} \approx 1.747\times10^{-8},
  \qquad
  x^* = \frac{R_0}{C^*} \approx 8.21\times10^{-10}.
  \label{P12:eq:physical_params_R0}
\end{equation}
The physical condensate torus is in the \emph{ultra-inverted} regime:
$C^*/R_0\approx1.22\times10^9$. This is not the thin-torus limit
but rather the opposite extreme --- a nearly-straight tube of radius
$C^*\approx21.28$ whose ends are identified at major radius
$R_0\approx1.75\times10^{-8}$ (all in natural units $\Lcond=1$).
\begin{remark}[Correction to earlier estimates]
\label{P12:rem:phys_correction}
An earlier version of this section estimated $R_0\approx0.041$ and
$x^*\approx0.002$. Those estimates were not derived from
the Theorem~\ref{P12:thm:E} constraint and are inconsistent with it:
they give $R_0 C^{*2}\approx18.6$, whereas the exact value is
$5/(\vp^{23}\pi^2)\approx7.91\times10^{-6}$, off by a factor
of $\sim2.4\times10^6$. Equations~\eqref{P12:eq:physical_params}
and~\eqref{P12:eq:physical_params_R0} give the correct values.
\end{remark}

\subsection{The consistency equation}

Combining Theorem~\ref{P12:thm:D} with $R_0=x\,C^*$
and the normalisation $\bst=300/\pi^2$, the reduced conjecture
\eqref{P12:eq:reduced} becomes:
\begin{equation}
  \boxed{
  \vp^{13/2}\cdot2^{2/3}\cdot(2\pi)^2\cdot x
  \sqrt{\frac{300}{\pi^2\,b^*(x)}}
  \left[\frac{16}{15}+\delta(x)\right]
  = 10,}
  \label{P12:eq:consistency}
\end{equation}
where $\delta(x)=C^{*4}I_{\mathrm{tor}}(x)/[R_0^2\cdot16/15]$
is the toroidal correction from Theorem~\ref{P12:thm:D}.
This is the \emph{profile normalisation consistency equation}:
a single transcendental equation for the one unknown $x=R_0/C^*$.
All coefficients are determined by $\vp$, $\pi$, and $b^*(x)$
through the exact Newton series of Theorem~\ref{P12:thm:C}.

\begin{remark}[Equivalence to the two-route proof]
\label{P12:rem:consistency_equiv}
Equation~\eqref{P12:eq:consistency} is the \emph{explicit form} of the
geometric constraint $R_0 C^{*2}=5/(\vp^{23}\pi^2)$ of Theorem~\ref{P12:thm:E}:
substituting $C^*=\sqrt{300/(\pi^2 b^*)}$ and $R_0=x C^*$ converts
\eqref{P12:eq:RC2} into~\eqref{P12:eq:consistency}.
The two-route proof of Section~\ref{P12:sec:E} establishes that any solution
to~\eqref{P12:eq:consistency} yields $\vp^9\sqrt{\Ja\Jfour}=10$;
the equilibrium axiom guarantees that the physical $x^*$ is that solution.
\end{remark}

\begin{corollary}[Closed form for $x^*=R_0/C^*$]
\label{P12:cor:xstar}
The physical aspect ratio $x^*=R_0/C^*$ is determined in closed form
by Theorem~\ref{P12:thm:E} and the CMB normalisation $\bst=300/\pi^2$:
\begin{equation}
  \boxed{x^* = \frac{\pi\,\tan^3(q^*)}{9600\sqrt{6}\,\vp^{23}},}
  \label{P12:eq:xstar_exact}
\end{equation}
where $q^*$ is the virial Newton series root of Theorem~\ref{P12:thm:C}.
Substituting the leading-order approximation $q^*\approx\pi/5$ and
using the pentagon identity $\tan(\pi/5)=\sqrt{\ms}/\vp$
(Theorem~\ref{P12:thm:B}) yields the leading-order closed form:
\begin{equation}
  \boxed{x^*_0 = \frac{\pi\,(3-\vp)^{3/2}}{9600\sqrt{6}\,\vp^{26}}
  \approx 7.996\times10^{-10},}
  \label{P12:eq:xstar_leading}
\end{equation}
which uses only $\vp$ and $\pi$.
The full series~\eqref{P12:eq:xstar_exact} achieves:
$x^*\approx8.209\times10^{-10}$ (exact),
with the leading form~\eqref{P12:eq:xstar_leading} accurate to $2.60\%$
and the one-step corrected form accurate to $0.010\%$.
\end{corollary}

\begin{proof}
From Theorem~\ref{P12:thm:E}: $R_0 C^{*2}=5/(\vp^{23}\pi^2)$.
With $C^*=\sqrt{\bst/b^*}=\sqrt{300/(\pi^2 b^*)}$ and $R_0=x^*\!C^*$:
\[
  x^*\!\cdot C^{*3} = \frac{5}{\vp^{23}\pi^2},
  \quad\Longrightarrow\quad
  x^* = \frac{5}{\vp^{23}\pi^2}\cdot\left(\frac{\pi^2 b^*}{300}\right)^{3/2}
      = \frac{5\pi\, b^{*3/2}}{\vp^{23}\cdot300^{3/2}}.
\]
Using $b^*=\tan^2(q^*)/8$ (Theorem~\ref{P12:thm:C}) and
$300^{3/2}=(8\cdot300)^{3/2}/8^{3/2}$:
\[
  x^* = \frac{5\pi\,\tan^3(q^*)/(8^{3/2})}{\vp^{23}\cdot300^{3/2}}
      = \frac{\pi\,\tan^3(q^*)}{\vp^{23}\cdot(8\cdot300)^{3/2}/5}.
\]
Since $(8\cdot300)^{3/2}=2400^{3/2}=2400\cdot20\sqrt{6}=48000\sqrt{6}$
and $48000/5=9600$, equation~\eqref{P12:eq:xstar_exact} follows.
The leading form~\eqref{P12:eq:xstar_leading} uses
$\tan(\pi/5)=\sqrt{\ms}/\vp$ (pentagon identity, Theorem~\ref{P12:thm:B}),
giving $\tan^3(\pi/5)=(\ms)^{3/2}/\vp^3=(3-\vp)^{3/2}/\vp^3$,
and $\vp^{23}\cdot\vp^3=\vp^{26}$.
\qed
\end{proof}

%──────────────────────────────────────────────────────────────
\section{Summary and the one-parameter chain}
\label{P12:sec:summary}
%──────────────────────────────────────────────────────────────

\begin{center}
\begin{longtblr}[
  caption={Results of Paper~XII.
    \textcolor{darkgreen}{Green} = proved.
    Input $J_4/J_a$ is recalled from Papers~II--III.},
  label={tab:results},
]{
  colspec={X[2.2,p] X[1.2,p] X[1,c]},
  rowhead=1,
}
\toprule
Statement & Source & Status \\
\midrule
$\Jfour/\Ja=2^{4/3}/\vp^5$ &
  Papers~II--III~\cite{Paper2,Paper3} &
  \textcolor{darkgreen}{Recalled} \\
$\rCMB=10\,\Lcond^4$ (arithmetic identity) &
  Thm~\ref{P12:thm:A} &
  \textcolor{darkgreen}{Proved} \\
Pentagon identities $\vp=2\cos(\pi/5)$, etc. &
  Thm~\ref{P12:thm:B} &
  \textcolor{darkgreen}{Proved} \\
Newton--icosahedral series for $b^*$ &
  Thm~\ref{P12:thm:C} &
  \textcolor{darkgreen}{Proved} \\
Exact toroidal integral $I_{\mathrm{tor}}=\frac{2}{15}\bsF(2,2;\frac{7}{2};-x^2)$ &
  Thm~\ref{P12:thm:D} &
  \textcolor{darkgreen}{Proved} \\
$I_{\mathrm{tor}}(1)=\pi-3$ identity &
  Thm~\ref{P12:thm:D} &
  \textcolor{darkgreen}{Proved} \\
Conjecture $\Leftrightarrow$ equilibrium $\Efb/V=\rCMB$ &
  Thm~\ref{P12:thm:E} &
  \textcolor{darkgreen}{Proved} \\
$\vp^9\sqrt{\Ja\Jfour}=10$ &
  Cor~\ref{P12:cor:proof} &
  \textcolor{darkgreen}{Proved} \\
$x^*=\pi\tan^3(q^*)/(9600\sqrt{6}\,\vp^{23})$ (exact) &
  Cor~\ref{P12:cor:xstar} &
  \textcolor{darkgreen}{Proved} \\
$x^*_0=\pi(3-\vp)^{3/2}/(9600\sqrt{6}\,\vp^{26})$ (leading order) &
  Cor~\ref{P12:cor:xstar} &
  \textcolor{darkgreen}{Proved} \\
\bottomrule
\end{longtblr}
\end{center}

With Corollary~\ref{P12:cor:proof} established, the one-parameter chain
is complete:
\begin{equation}
  \boxed{
  \TCMB
  \;\xrightarrow{\;\Lcond\;}\;
  \rCMB=10
  \;\xrightarrow{\;\Efb/V=\rCMB\;}\;
  \vp^9\sqrt{\Ja\Jfour}=10
  \;\xrightarrow{\;\text{Paper~I}\;}\;
  m_e = \Lcond\,\vp^{20}\,e^{4\pi-3/400-\alpha/\pi}.
  }
  \label{P12:eq:chain}
\end{equation}
The electron mass is derivable from the CMB temperature $\TCMB$ alone,
with no free parameters. The three terms in the exponent are:
$4\pi$ (Chern--Simons spoke amplitude, exact); $-3/400$ (WZW modular
$T$-matrix phase $T_{1/2}/Q$, exact); and $-\alpha/\pi$ (one-loop QED
pole-mass to $\overline{\mathrm{MS}}$ scheme conversion, Paper~IV~\cite{Paper4}).
The residual after all three corrections is $0.013\%$,
the irreducible $O(R_0^{-2})$ thick-torus geometric correction.

\paragraph{What the proof uses.}
Four ingredients, all independent of numerical fitting:
\begin{enumerate}[label=(\arabic*),noitemsep]
\item $\rCMB=10\,\Lcond^4$ (Theorem~\ref{P12:thm:A}: pure arithmetic).
\item $\Jfour/\Ja=2^{4/3}/\vp^5$ (Papers~II--III: topological ratio).
\item $\Efb=\Ja\Jfour/\vp^5$ (condensate energy from the feedback functional).
\item $\Efb/V_{\mathrm{torus}}=\rCMB$ (thermodynamic equilibrium axiom).
\end{enumerate}
Items (1)--(3) are proved results; item (4) is the physical
equilibrium condition that selects the physical saddle.
The $\pi-3$ identity (Theorem~D) and the Newton series (Theorem~C)
provide the exact geometric description of the saddle and the
convergent expansion for $b^*$, but are not needed for the
conjecture proof itself.

\paragraph{Relation to the series.}
Paper~XII closes the last logical gap in the main condensate programme:
the profile normalisation conjecture and the one-parameter chain.
The chain $\TCMB\to m_e$ is now a theorem: the single input
$\TCMB=2.7255\,\mathrm{K}$ determines the condensate scale $\Lcond$,
which via the equilibrium condition forces $\vp^9\sqrt{\Ja\Jfour}=10$,
which via Paper~I's Hopf spoke gives $m_e/\Lcond=\vp^{20}e^{4\pi-3/400}$.
The subsequent papers (IX, X, XI) then extend the framework to
quantum mechanics, Bell inequalities, and atomic physics ---
all following from the single input $\TCMB$
established here.

%──────────────────────────────────────────────────────────────
\paragraph{What remains open.}
Paper~XII resolves the conjecture and closes the one-parameter chain,
but two related problems remain open:
\begin{enumerate}[label=(\arabic*),noitemsep]
\item \emph{$O(R_0^{-2})$ corrections to physical observables.}
  The closed form $x^*=\pi\tan^3(q^*)/(9600\sqrt{6}\,\vp^{23})$
  (Corollary~\ref{P12:cor:xstar}) determines the torus geometry exactly.
  The remaining open problem is to evaluate the
  $O(R_0^{-2})$ corrections to physical observables ($m_e$, $v_{\rm EW}$)
  at the exact geometry, which requires the full 2D
  Euler--Lagrange profile $f^*(\varrho,\theta)$ beyond the
  Battye--Sutcliffe ansatz.
\item \emph{$O(R_0^{-2})$ corrections to physical observables.}
  The residuals $-0.013\%$ in $m_e$ (Paper~IV) and $0.69\%$ in
  $v_{\rm EW}$ (Paper~I) are both $O(R_0^{-2})$ corrections
  from the thick-torus geometry at the physical $x^*\approx8.21\times10^{-10}$,
  not captured by the Battye--Sutcliffe ansatz.
  Closing them requires the exact 2D Euler--Lagrange profile.
\end{enumerate}
Both items are subleading corrections to quantities already established;
they do not affect the validity of the conjecture proof or the
one-parameter chain.

\bibliographystyle{ieeetr}
\begin{thebibliography}{9}
\bibitem{Paper1}
F.~Manfredi,
\textit{The Density-Feedback Faddeev--Niemi Hopfion:
  Exact Algebraic Predictions for Standard-Model Parameters},
Zenodo (2026).
\href{https://doi.org/10.5281/zenodo.19342027}{doi:10.5281/zenodo.19342027}

\bibitem{Paper2}
F.~Manfredi,
\textit{BPS Structure and Fixed-Point Theorem for the
  Density-Feedback Faddeev--Niemi Hopfion},
Zenodo (2026).
\href{https://doi.org/10.5281/zenodo.19363491}{doi:10.5281/zenodo.19363491}

\bibitem{Paper3}
F.~Manfredi,
\textit{Two Constants from One Knot: The Fine Structure Constant and
  the Linking Scale as Exact Predictions of the Icosahedral WZW Condensate},
Zenodo (2026).
\href{https://doi.org/10.5281/zenodo.19478629}{doi:10.5281/zenodo.19478629}

\bibitem{Paper4}
F.~Manfredi,
\textit{The Hopf Spoke, WZW Fermion Mass Renormalization,
  and Three- and Four-Term Formulas for the Fine Structure Constant},
Preprint (2026).
\href{https://doi.org/10.5281/zenodo.19504729}{doi:10.5281/zenodo.19504729}

\bibitem{PaperFC}
F.~Manfredi,
\textit{Foundational Companion: The Pentagon Theorem and Three-Route
Proof of $k=3$},
Preprint (2026).
\href{https://doi.org/10.5281/zenodo.20173651}{doi:10.5281/zenodo.20173651}

\bibitem{Battye1998}
R.~A. Battye and P.~M. Sutcliffe,
\textit{Knots as stable soliton solutions in a three-dimensional classical
field theory},
Phys.\ Rev.\ Lett.\ \textbf{81}, 4798 (1998).
\href{https://doi.org/10.1103/PhysRevLett.81.4798}{doi:10.1103/PhysRevLett.81.4798}

\bibitem{AS65}
M.~Abramowitz and I.~A. Stegun,
\textit{Handbook of Mathematical Functions},
Dover, New York (1965).
ISBN 978-0-486-61272-0.
\end{thebibliography}
\end{document}